\documentclass[11pt]{article}
\usepackage[margin=2.5cm]{geometry}
\usepackage{amsmath}
\usepackage{parskip}

\title{T5A2}
\author{}
\date{}

\begin{document}
\maketitle

The data are $f_i = f(x_i)$ for $i = 0, 1, \dots, n$. I approximate them with the polynomial
\[
p_m(x) = \sum_{k=0}^{m} c_k x^k,
\]
and I look for the coefficients $c_0, c_1, \dots, c_m$ that make
\[
\chi^2 = \sum_{i=0}^{n} [p_m(x_i) - f_i]^2
\]
as small as possible.

1) Analytical formulas for the coefficients

At the minimum of $\chi^2$, the derivative to every coefficient is zero:
\[
\frac{\partial \chi^2}{\partial c_l} = 0, \qquad l = 0, 1, \dots, m.
\]
In $p_m(x_i) = c_0 + c_1 x_i + \dots + c_m x_i^m$ only the term $c_l x_i^l$ depends on $c_l$, so $\partial p_m(x_i)/\partial c_l = x_i^l$. With the chain rule,
\[
\frac{\partial \chi^2}{\partial c_l} = \sum_{i=0}^{n} 2[p_m(x_i) - f_i] x_i^l = 0.
\]
I divide by 2 and put in $p_m(x_i) = \sum_k c_k x_i^k$:
\[
\sum_{i=0}^{n} \sum_{k=0}^{m} c_k x_i^{k} x_i^{l} = \sum_{i=0}^{n} f_i x_i^l.
\]
I swap the two sums, and I take $c_k$ out of the sum over $i$ because it does not depend on $i$:
\[
\sum_{k=0}^{m} \left(\sum_{i=0}^{n} x_i^{l+k}\right) c_k = \sum_{i=0}^{n} f_i x_i^l, \qquad l = 0, 1, \dots, m.
\]
These sums only depend on the data, so I call them
\[
s_j = \sum_{i=0}^{n} x_i^{j} \quad (j = 0, 1, \dots, 2m), \qquad t_l = \sum_{i=0}^{n} f_i x_i^{l} \quad (l = 0, 1, \dots, m).
\]
I need $s_j$ up to $j = 2m$, because $l + k$ goes up to $2m$. Note that $s_0 = n+1$, because $x_i^0 = 1$. Now the equations are
\[
\sum_{k=0}^{m} s_{l+k} c_k = t_l, \qquad l = 0, 1, \dots, m.
\]
These are $m+1$ linear equations for the $m+1$ unknowns $c_k$. In matrix form, $S c = t$:
\[
\begin{pmatrix}
s_0 & s_1 & \cdots & s_m \\
s_1 & s_2 & \cdots & s_{m+1} \\
\vdots & \vdots & \ddots & \vdots \\
s_m & s_{m+1} & \cdots & s_{2m}
\end{pmatrix}
\begin{pmatrix} c_0 \\ c_1 \\ \vdots \\ c_m \end{pmatrix}
=
\begin{pmatrix} t_0 \\ t_1 \\ \vdots \\ t_m \end{pmatrix},
\]
so $S_{lk} = s_{l+k}$. If $\det S \neq 0$, Cramer's rule gives every coefficient:
\[
c_k = \frac{\det S^{(k)}}{\det S}, \qquad k = 0, 1, \dots, m,
\]
where $S^{(k)}$ is $S$ with column $k$ replaced by $t$.

Now I check that $\det S \neq 0$ when at least $m+1$ of the $x_i$ are different (so $n \geq m$). I take any vector $v = (v_0, v_1, \dots, v_m) \neq 0$ and the polynomial $q(x) = \sum_{k=0}^{m} v_k x^k$. With $s_{l+k} = \sum_i x_i^l x_i^k$,
\[
v^T S v = \sum_{l=0}^{m} \sum_{k=0}^{m} v_l s_{l+k} v_k
= \sum_{i=0}^{n} \left(\sum_{l=0}^{m} v_l x_i^l\right)\left(\sum_{k=0}^{m} v_k x_i^k\right)
= \sum_{i=0}^{n} q(x_i)^2.
\]
$q$ is not the zero polynomial and its degree is at most $m$, so it has at most $m$ roots. So $q(x_i) \neq 0$ for at least one of the $m+1$ different $x_i$, and $v^T S v > 0$.
If $\det S = 0$, there would be a $v \neq 0$ with $S v = 0$, and then $v^T S v = 0$. That is not possible, so $\det S \neq 0$.

I also check that this solution is a minimum of $\chi^2$. I change the coefficients to $c_k + \delta_k$ with $\delta \neq 0$. Then $p_m(x_i)$ changes by $q(x_i) = \sum_k \delta_k x_i^k$, and with $(a+b)^2 = a^2 + 2ab + b^2$:
\[
\begin{aligned}
\chi^2(c + \delta) &= \sum_{i=0}^{n} [p_m(x_i) - f_i + q(x_i)]^2 \\
&= \sum_{i=0}^{n} [p_m(x_i) - f_i]^2 + \sum_{l=0}^{m} \delta_l \left(\sum_{i=0}^{n} 2[p_m(x_i) - f_i] x_i^l\right) + \sum_{i=0}^{n} q(x_i)^2 \\
&= \chi^2(c) + \sum_{l=0}^{m} \delta_l \frac{\partial \chi^2}{\partial c_l} + \delta^T S \delta.
\end{aligned}
\]
The middle term is zero, because every $\partial \chi^2 / \partial c_l = 0$ at the solution. The last term is $\delta^T S \delta > 0$, the same as $v^T S v$ above. So $\chi^2(c + \delta) > \chi^2(c)$ for every $\delta \neq 0$, and the solution is the minimum.

2) Main steps of the numerical algorithm

Numerically I do not use Cramer's rule. It needs $m+2$ determinants, and one determinant costs about as much as one Gaussian elimination. Solving $S c = t$ directly needs only one elimination.

Step 1: Input the data $x_i$, $f_i$ for $i = 0, 1, \dots, n$ and the order $m$. At least $m+1$ of the $x_i$ must be different, so that $\det S \neq 0$ (see part 1).

Step 2: Compute the sums $s_j$ for $j = 0, 1, \dots, 2m$ and $t_l$ for $l = 0, 1, \dots, m$ in one loop over the data points. I get the powers by multiplying with $x_i$ every time, $x_i^{j+1} = x_i^j \cdot x_i$.

Step 3: Fill the $(m+1) \times (m+1)$ matrix $S$ with $S_{lk} = s_{l+k}$, and the vector $t$.

Step 4: Solve $S c = t$ with Gaussian elimination with partial pivoting (section 5.3 in the book). For each column $j$, I swap row $j$ with the row (from $j$ down) that has the largest $|S_{rj}|$, and I swap the same two elements of $t$. This way I never divide by a zero or very small pivot. Then I subtract $S_{rj}/S_{jj}$ times row $j$ from every row $r$ below it (also in $t$), so everything below the diagonal becomes zero. After that I find the coefficients with back substitution, from $c_m$ to $c_0$:
\[
c_l = \frac{1}{S_{ll}}\left(t_l - \sum_{k=l+1}^{m} S_{lk} c_k\right),
\]
with $S$ and $t$ as they are after the elimination.

Step 5: Output $c_0, c_1, \dots, c_m$.

3) Subprogram

The subprogram for steps 1 to 5, and the Gaussian elimination it uses:

\begin{verbatim}
subprogram leastSquares(x[0..n], f[0..n], m) returns c[0..m]
    # Step 1: the input is the data x_i, f_i and the order m
    # Step 2: the sums s_j and t_l
    s[0..2m] = 0
    t[0..m] = 0
    for i = 0 to n #loop over the data points
        xp = 1 #xp = x_i^j, starts with x_i^0 = 1
        for j = 0 to 2m
            s[j] = s[j] + xp
            if j <= m: t[j] = t[j] + f[i]*xp #t only goes up to m
            xp = xp*x[i] #next power of x_i
    # Step 3: fill the matrix S with S_lk = s_(l+k)
    for l = 0 to m
        for k = 0 to m
            S[l][k] = s[l+k]
    # Step 4: solve S c = t
    c = gauss(S, t)
    # Step 5: output the coefficients
    return c

subprogram gauss(S[0..m][0..m], t[0..m]) returns c[0..m]
    # elimination with partial pivoting, column by column
    for j = 0 to m-1
        # largest pivot, so I never divide by zero or a very small number
        piv = the row r >= j with the largest |S[r][j]|
        swap rows j and piv of S, and t[j] and t[piv]
        for r = j+1 to m
            fac = S[r][j]/S[j][j] #chosen so that S[r][j] becomes 0
            for k = j to m
                S[r][k] = S[r][k] - fac*S[j][k]
            t[r] = t[r] - fac*t[j] #same row operation on t
    # S is upper triangular now, so back substitution from the last row up
    for l = m down to 0
        c[l] = (t[l] - sum of S[l][k]*c[k] for k = l+1 to m)/S[l][l]
    return c
\end{verbatim}

\end{document}
